% TOP_PROBLEM: {"id": 30006679, "title": "NP versus coNP Problem", "category_id": 15, "category": {"id": 15, "name": "computer_science", "display_name": "Computer Science", "description": "Computational complexity, algorithms, and theoretical CS.", "slug": "computer-science", "order_index": 15, "created_at": "2026-07-31T15:26:25.671Z"}, "status": "open"}
% FIRST_POSTED: 2026-09-27
% !TeX program = lualatex
% ATTEMPT_STATUS: unresolved
\documentclass[11pt]{article}
\usepackage[margin=25mm]{geometry}
\usepackage{fontspec}
\setmainfont{Latin Modern Roman}
\usepackage{amsmath,amssymb,mathtools}
\usepackage{unicode-math}
\setmathfont{Latin Modern Math}
\usepackage{xurl,hyperref}
\hypersetup{colorlinks=true,urlcolor=blue}
\setcounter{secnumdepth}{0}
\setlength{\parindent}{0pt}
\setlength{\parskip}{5pt}
\setlength{\emergencystretch}{3em}
\begin{document}
\section*{\texorpdfstring{NP versus coNP Problem}{NP versus coNP Problem}}\label{30006679}
\subsection{Definitions and mathematical statement}
For $L\subseteq\{0,1\}^*$ write $\overline L=\{0,1\}^*\setminus L$ and set $\mathsf{coNP}=\{\overline L:L\in\mathsf{NP}\}$. The question is $\mathsf{NP}=\mathsf{coNP}$. A Cook--Reckhow proof system is a polynomial-time computable surjection $P:\{0,1\}^*\to\mathsf{TAUT}$, where $\mathsf{TAUT}$ is the set of propositional tautologies. Polynomial boundedness means
\[
 \exists\text{ polynomial }q\ \forall\varphi\in\mathsf{TAUT}\ \exists\pi:
 P(\pi)=\varphi,\quad|\pi|\le q(|\varphi|).
\]
The equivalent proof-system question permits choosing the system; it is not restricted to a familiar calculus.
\par\smallskip\textit{Formulation and concept references:} \href{https://www.scottaaronson.com/papers/pnp.pdf}{[S1]}.

\subsection{Short English statement}
Does efficient certification of yes-answers always come with efficient certification of no-answers?
\subsection{Research attempt}
\paragraph{Scope, source check, and plan.}
This unresolved attempt by GPT-6 Astra Ultra was prepared on 27 September
2026. The process was to check the proof-system formulation, examine a
concrete certificate-compression strategy, and construct adversarial examples
for conclusions drawn from representation size. No separation or collapse of
complexity classes is claimed.

Cook--Reckhow [E1], Propositions 1.1 and 1.4, give the equivalence between
$\mathsf{NP}=\mathsf{coNP}$ and existence of a polynomially bounded proof
system. Their definition ranges over all polynomial-time surjections onto
tautologies. The catalogue's [S3] explicitly restricts its lower-bound target
to tree-like semantic systems with bounded line size. The other catalogue
research entries are retained as references; their titles are not accepted
as proofs of a class equality or separation.

The selected route is to test whether short unsatisfiability certificates can
be obtained by eliminating auxiliary variables from a verifier. Parity gives
an exact calculation: elimination can require exponentially many clauses,
while an elementary algebraic certificate remains short. A separate explicit
proof-system construction then tests whether a lower bound in one verifier
can possibly quantify over all verifiers.

\paragraph{The universal quantifier over proof systems.}
If a proof system $P$ has proofs bounded by a polynomial $q$, a verifier for
$\mathsf{TAUT}$ guesses $\pi$ of length at most $q(|\varphi|)$ and checks
$P(\pi)=\varphi$. Conversely, if $\mathsf{TAUT}$ has a polynomially bounded
NP verifier $V$, define $P(\varphi,w)=\varphi$ when $V(\varphi,w)$ accepts,
and output a fixed tautology $\tau_0$ otherwise. This is a polynomial-time
surjection with polynomially bounded proofs. Together with coNP-completeness
of tautology, this proves the equivalence used here [E1]. Consequently a
negative answer requires failure of polynomial boundedness for \emph{every}
such $P$. A positive answer needs just one sound, complete, polynomial-time
verifier with a uniform polynomial proof bound.

\paragraph{Lemma 13.1: exact clause complexity of parity.}
For $n\ge1$ let
\[
 E_n=\{x\in\{0,1\}^n:x_1+\cdots+x_n=0\pmod2\}.
\]
Every CNF on these $n$ variables alone whose satisfying set is exactly $E_n$
contains at least $2^{n-1}$ non-tautological clauses. This bound is attained.

\emph{Proof.} Remove repeated literals and tautological clauses. Every
remaining clause must be true on all even-parity inputs. If a clause omits
$x_i$, assign the variables it mentions so as to falsify every literal.
Choose the other unspecified variables arbitrarily except $x_i$, and choose
$x_i$ last to make the total parity even. This gives a satisfying input of
$E_n$ falsifying the clause, a contradiction. Thus every such clause mentions
all $n$ variables. With no repeated or complementary literals, it is falsified
by exactly one assignment; that assignment must have odd parity. All
$2^{n-1}$ odd assignments must be excluded, giving the lower bound.
Conversely, for each odd assignment take the width-$n$ clause falsified
precisely by it. Their conjunction defines $E_n$.\hfill$\square$

The result also holds for odd parity, by exchanging the final parity bit.
It is a lower bound on an exact visible-variable representation, with no
hypothesis about how that representation is generated. It therefore falsifies
any proposed universal procedure that projects an arbitrary short CNF onto
its visible variables and always returns a polynomial-size equivalent CNF
without new variables.

\paragraph{A short extended representation, checked gate by gate.}
Introduce prefix variables $p_0,\ldots,p_n$, require $p_0=0$, and impose
\[
                   p_i=p_{i-1}+x_i\pmod2\quad(1\le i\le n).
\]
Each relation $c=a+b\pmod2$ is given by the four clauses
\[
 (a\vee b\vee\neg c),\quad
 (a\vee\neg b\vee c),\quad
 (\neg a\vee b\vee c),\quad
 (\neg a\vee\neg b\vee\neg c).
\]
They exclude exactly the four triples violating the XOR equation. The unit
clauses $\neg p_0$ and $\neg p_n$ complete a CNF with $4n+2$ clauses.
For every input $x$, the gate equations determine a unique sequence $p_i$;
its final value is the parity of $x$. Thus existentially projecting the
$p$ variables gives exactly $E_n$. With binary variable indices, the encoded
length is $O(n\log(n+2))$.

The exponential lower bound is consequently not a statement that parity
requires long descriptions, or that checking parity is difficult. It records
the cost of requiring a particular description after forgetting the internal
prefix values. This distinction matters in proof search: a short proof can
retain useful intermediate variables instead of projecting them away.

\paragraph{Lemma 13.2: short certificates for affine inconsistency.}
Let $A\in\mathbb F_2^{m\times r}$ and $b\in\mathbb F_2^m$. The system
$Az=b$ is inconsistent if and only if there is
$\lambda\in\mathbb F_2^m$ with
\begin{equation}\label{a13:dual}
                  \lambda^{\mathsf T}A=0,
                  \qquad\lambda^{\mathsf T}b=1.
\end{equation}
This is an $m$-bit certificate checked in $O(mr+m)$ field operations.

\emph{Proof.} Such a certificate would transform a hypothetical solution
into $0=1$, so it proves inconsistency. Conversely apply Gaussian row
elimination to the augmented matrix $(A\mid b)$, tracking each row operation
on an additional identity matrix $I_m$. If elimination never produces a zero
coefficient row with right side $1$, assign free variables arbitrarily and
solve the pivot equations backwards, obtaining a solution. In the
inconsistent case there is therefore a row $(0\mid1)$; the corresponding
tracked row of the identity matrix is $\lambda^{\mathsf T}$. Row operations
preserve the stated linear-combination interpretation. All matrices and all
operations have polynomial size.\hfill$\square$

Apply this directly to two prefix encodings on the same $x$ variables, with
$p_0=q_0=0$, respective equations
\[
 p_i+p_{i-1}+x_i=0,\qquad q_i+q_{i-1}+x_i=0,
\]
and the incompatible requirements $p_n=0$, $q_n=1$. Sum every displayed
prefix equation and all four boundary equations. Each $x_i$ occurs twice;
every internal $p_i,q_i$ occurs twice; the boundary occurrences also cancel.
The right side is $1$. Thus the all-ones row-combination vector is already
an explicit certificate for $0=1$. It has $2n+4$ bits.

A procedure that first expands both projected parity conditions into CNF
would generate exponentially many clauses before noticing their conflict.
The affine certificate avoids that expansion. This demonstrates why the
representation lower bound in Lemma 13.1 cannot by itself lower-bound proof
length across all proof systems.

\paragraph{Where the algebraic upper-bound proposal breaks.}
Boolean circuits do not consist only of XOR relations. Over $\mathbb F_2$,
an AND gate has equation $z+xy=0$. Treating the monomial $xy$ as an independent
linear variable discards its relation to $x$ and $y$. There is an immediate
counterexample: the equations
\[
                    x=0,\quad y=0,\quad z=1,\quad z=xy
\]
are inconsistent, but replacing $xy$ by a free variable $u$ gives the
consistent affine system $x=0,y=0,z=1,z=u$. Therefore affine inconsistency
certificates remain sound for the affine systems to which they apply, but
this naive linearization is incomplete for Boolean gate systems. Restoring
$u=xy$ restores the nonlinear constraint. No polynomial bound for certificates
handling every such nonlinear system is proved here.

\paragraph{An explicit stress test for proof-system lower bounds.}
There is also a simple way to make the quantifier failure concrete. Define
a truth-table proof system $T$: a well-formed proof contains a formula
$\varphi$ and the evaluations on all assignments in a fixed order; $T$
checks the table and outputs $\varphi$ if all entries are $1$. Otherwise it
outputs $\tau_0$. The checker first verifies that the supplied table length
is exactly $2^v$, where $v$ is the number of distinct variables. It performs
the exponential-in-$v$ scan only when that many bits were supplied. Hence
its running time is polynomial in the \emph{proof} length, as required, and
it is a surjection onto tautologies.

Let $\tau_n=\bigwedge_{i=1}^n(x_i\vee\neg x_i)$, choosing $n$ so that this
formula is not the default $\tau_0$. Its formula length is
$O(n\log(n+2))$, but every $T$-proof has at least $2^n$ table entries. Now
extend $T$ with a tagged input containing any formula of the syntactic form
$\tau_n$, which is output immediately. Recognition and output take
polynomial time, and every added output is a tautology. The resulting system
is still sound and surjective, and the same family has proofs of length
$O(n\log(n+2))$. Thus a proved exponential lower bound for a particular
system and family can coexist with short proofs in another permitted system.
This is an illustrative construction, not a claim that any familiar strong
system simulates all others.

\paragraph{Verification and first remaining gap.}
The companion checks enumerate all normalized clauses on up to six variables,
confirm that every non-tautological clause valid on even parity has full
width, and count exactly $2^{n-1}$ necessary excluded odd assignments.
They test all eight XOR triples and the four displayed gate clauses.
They also enumerate all $2\times3$ affine systems, comparing the existence
of a solution with the absence of a certificate \eqref{a13:dual}. These
checks support the finite identities, while the proofs cover arbitrary sizes.

The first missing lemma for the proposed upper-bound route is a sound,
complete polynomial-time verifier whose certificates for \emph{all}
unsatisfiable Boolean gate systems have length polynomial in the original
encoding. The affine lemma proves this only for a tractable subclass. For
a lower-bound route, the missing quantifier is a bound against every
Cook--Reckhow system, including systems with additional certificate types.
Next work must address nonlinear constraints, retained auxiliary variables,
and this universal quantifier together; more accurate parity clause counts
will not bridge the gap.

\paragraph{Outcome.}
\textbf{UNRESOLVED.} The attempt proves an exact parity representation bound,
constructs polynomial certificates for an affine contradiction family, and
exhibits a proof-system extension defeating a particular exponential lower
bound. These elementary results are presented without novelty claims. They
identify two invalid inferences, but establish neither
$\mathsf{NP}=\mathsf{coNP}$ nor $\mathsf{NP}\ne\mathsf{coNP}$.

\subsection{Sources}
\begin{itemize}
\item[S1]Scott Aaronson, "P =? NP," in Open Problems in Mathematics (J. F. Nash Jr., M. Th. Rassias, eds.), Springer, 2016, pp. 1-122.
\url{https://www.scottaaronson.com/papers/pnp.pdf}.
\textit{Formulation source.}
\item[S2]Proof Complexity and Feasible Interpolation — Amirhossein Akbar Tabatabai (2025).
\url{https://arxiv.org/html/2505.03002v1}.
\textit{Further reference; not a proof endorsement.}
\item[S3]Superpolynomial Length Lower Bounds for Tree-Like Semantic Proof Systems with Bounded Line Size — de Rezende, Engström, Ghannane and Risse (2026).
\url{https://eccc.weizmann.ac.il/report/2026/078/}.
\textit{Further reference; not a proof endorsement.}
\item[S4]Simulating Polynomial-Time Nondeterministic Turing Machines via Nondeterministic Turing Machines — Tianrong Lin, v28.
\url{https://arxiv.org/abs/2406.10476}.
\textit{Further reference; not a proof endorsement.}
\item[S5]A Critique of Lin's On NP versus coNP and Frege Systems — DeJesse, Lyudovyk, Pai and Reidy (2025).
\url{https://arxiv.org/abs/2505.05658}.
\textit{Further reference; not a proof endorsement.}
\item[S6]Proofs of NP = coNP = PSPACE: Current upgrade — Lev Gordeev and Edward Hermann Haeusler, v3.
\url{https://arxiv.org/abs/2311.17939}.
\textit{Further reference; not a proof endorsement.}
\item[S7]A simplified lower bound for implicational logic — Emil Jeřábek (September 2024 version; published 2025).
\url{https://arxiv.org/html/2303.15090v3}.
\textit{Further reference; not a proof endorsement.}
\item[E1]Stephen A. Cook and Robert A. Reckhow,
\emph{The Relative Efficiency of Propositional Proof Systems},
Journal of Symbolic Logic 44(1) (1979), pp. 36--50, especially
Propositions 1.1 and 1.4 and Definition 1.3.
\url{https://www.cs.toronto.edu/~sacook/homepage/cook_reckhow.pdf}.
\url{https://doi.org/10.2307/2273702}.
\textit{Primary source for the proof-system equivalence and definition;
consulted 27 September 2026.}
\end{itemize}
\end{document}
